\chapter{Conclusion}\label{ch:conclusion}

This dissertation began with a practical problem: systems engineers increasingly want to use and stand to benefit from \acp{llm} to accelerate work across the digital thread, yet there is no established evaluation method for deciding where such tools are trustworthy and where they introduce unacceptable risk beyond relying on general-purpose benchmarks. Most evaluation practice relies on multiple-choice benchmarking because it is scalable and straightforward to score. However, many systems engineering activities are not selection problems. They require generating and justifying free-text outputs---rationale, explanations, trade studies, and narrative evidence---that must remain coherent with the digital thread of artifacts. In these contexts, multiple-choice results can inflate perceived readiness by measuring constrained recognition rather than generative adequacy. Open-style evaluation can better test free-text capability, but brings practical burdens in dataset engineering, scoring design including LLM-as-a-judge protocols, and computational cost.

The research performed in this dissertation re-frames model selection as an evidence-based assessment of suitability about intended use. The central claim is that \textbackslash{}ac\{mcq\} and \textbackslash{}ac\{osq\} should be treated as complementary evidence rather than functional equivalents: the two modalities impose different response constraints and therefore reveal different limitations even under aligned knowledge targets. To make this claim testable, SysEngBench is extended into a paired, multi-format evaluation instrument with controlled \ac{mcq} variants and matched \acp{osq}. This instrument enables three analyses central to qualification: (i) robustness diagnostics under \ac{mcq} structural perturbations, (ii) cross-modality evidence comparisons that quantify modality separation by task category, and (iii) token-efficiency analysis that makes cost and feasibility explicit. These contributions are delivered through a reproducible end-to-end pipeline spanning dataset construction, standardized inference, rubric-based scoring, and statistical analysis.

The remainder of this chapter translates these findings into practical guidance for systems engineers, summarizes the dissertation’s contributions, and closes with final reflections on evidence-based adoption of \acp{llm} in the digital thread.


\section{Practical Applications for Systems Engineering and Insights for LLM Deployment}\label{sec:practical_applications}

This dissertation’s central deployment takeaway is that benchmark evidence is most actionable when translated into \emph{decision rules}. Rather than selecting a single ``best'' model and treating it as uniformly reliable, practitioners can operationalize the findings through three linked practices: (i) a task-to-modality routing map that determines what evidence is required for a task class, (ii) a two-stage qualification workflow that uses \ac{mcq} for scalable screening and \ac{osq} for deeper generative qualification when warranted, and (iii) cost-aware selection and budgeting using token-efficiency criteria within the set of qualified candidates. Collectively, these practices turn the dissertation’s qualification framing into an actionable SE approach: they align evaluation evidence with the intended use claim and restrict model choice to options that satisfy the practitioner's risk posture for adoption and operational constraints. 

% ---------------------------------------------------------
\subsection{Model Selection Guidance by Task Type}

From a practical standpoint, model selection should be guided not only by average benchmark performance, but by how well a model’s measured strengths align with the target task’s evidence requirements. This work supports making that alignment explicit through a \emph{task-to-modality routing map}: a category-driven rule set that specifies the minimum defensible qualification modality for each task class. The routing map treats evaluation modality as a measurement instrument choice. In categories where cross-modality separation is high (e.g., where \ac{mcq} accuracy materially overstates \ac{osq} generative adequacy), \ac{mcq}-only evidence is insufficient to justify deployment for artifact-impacting work. In contrast, categories with smaller observed separation may permit \ac{mcq}-dominant screening, but still benefit from targeted \ac{osq} qualification where higher confidence is required. In this way, modality selection becomes a first-order part of system safety and governance: the organization decides \emph{what kind of evidence is required} before deciding \emph{which model to use}.

Operationally, the routing map can be implemented as a simple policy table that maps task families (or \ac{incose} category groupings) to an evidence requirement (e.g., \ac{mcq}-screening only versus dual-modality qualification). The benefit is that it counters a common shortcut in practice: treating multiple-choice benchmark performance as a stand-in for readiness across the full range of systems engineering work products. By routing tasks to modality-appropriate evidence, practitioners gain a principled way to choose which model to use for a given task class, prioritizing candidates that demonstrate strong generative quality, completeness, and justified outputs where those attributes matter. 

% ---------------------------------------------------------
\subsection{Benchmark-Informed SE Workflow Integration}

Benchmark results provide the evidence needed to \textbf{select} \textbackslash{}acp\{llm\} for integration into systems engineering workflows as \textbackslash{}emph\{role-assigned copilots\} rather than general-purpose assistants. Consistent with the dissertation’s qualification framing, the findings support \emph{evidence fusion}: \ac{mcq} and \ac{osq} provide overlapping but non-identical signals, and each is informative for selecting models that fit specific workflow roles. \ac{mcq}-style evidence offers scalable measurement of constrained-response competence and can be augmented with robustness diagnostics to assess stability under structural perturbations. \ac{osq}-style evidence directly probes generative adequacy for free-text synthesis, rationale generation, and narrative artifacts that may enter the digital thread.

These techniques need not be applied in a single fixed sequence. Instead, they can be composed based on what must be decided: narrowing the candidate set, differentiating between closely performing models, or validating readiness for tasks where generative adequacy is the primary success criterion. In practice, this supports modality-aware role assignment: models that demonstrate consistent performance under constrained selection tasks can be used for conformance checks and structured validation prompts, while models that score strongly on \ac{osq} can be used for drafting explanations, justification, and other narrative work products. 

This perspective generalizes naturally to \ac{mbse}. As models and documents become repository-managed and increasingly text-addressable, copilot design can follow established software engineering patterns: bounded change sets, schema-constrained outputs, automated checks, and human review at defined decision points. Benchmark evidence becomes the justification for \emph{where} automation could be leveraged successfully and \emph{how} it is bounded, rather than a claim that the model is uniformly reliable.

% ---------------------------------------------------------
\subsection{Budgeting and Token Constraints}

The tokenomics results translate directly into deployment guidance by making feasibility and cost explicit. Response verbosity varies widely across models, and longer responses do not reliably produce higher judged quality. This supports treating token usage as an operational budget variable and selecting models through a cost-aware rule that is consistent with the qualification framing: first qualify for a given task, then optimize. Specifically, models are filtered for acceptable \ac{osq} performance (e.g., at or above a ``Good'' or ``Pass'' threshold), and then compared on efficiency for that task using the quality--token trade space visualized in Figure~\ref{fig:pareto_quality_vs_token_usage}. Under comparable judged performance, the most cost-effective choice is the model that meets the required quality threshold using fewer median tokens per response.

Practically, this enables more control over cost. Practitioners can translate tokenomics evidence into cost-conscious model deployments by selecting the more efficient model(s) for a given task. The resulting guidance is straightforward: \emph{select the lowest-token model among those meeting the required quality threshold}, and enforce token and format constraints as part of the workflow’s configuration-controlled operating procedure.


\section{Summary of Contributions}\label{sec:summary_contributions}

Chapter~\ref{sec:research_contributions} defined four contributions. This chapter summarizes those contributions in the same order, emphasizing what was delivered and why each contribution matters in light of the dissertation’s empirical results and the deployment-oriented practical applications derived from Chapters~\ref{ch:introduction}--\ref{ch:results}. Collectively, these contributions formalize model selection for systems engineers as a qualification problem by translating measurement into decision rules. These rules include task-to-modality routing, modality-aware evidence composition for workflow integration, and cost-aware model selection under token budgets. 

Chapter~\ref{ch:introduction} provided a compact problem-to-contributions overview in Table~\ref{tab:rq_obj_contrib_mapping}. For readers who want a more detailed audit trail and navigation aid, Appendix~\ref{app:traceability_map} provides a traceability map linking research questions to objectives, methods, results, implications, and the four contributions detailed below.

\subsection{Extension of SysEngBench with Multi-Format Evaluation Questions}\label{subsec:ch6_contrib1}

This dissertation extended SysEngBench beyond a single-format \ac{mcq} benchmark by introducing a paired, multi-format evaluation instrument consisting of (i) controlled \ac{mcq} variants and (ii) \acp{osq} derived from the same underlying systems engineering knowledge targets. Rather than treating formats as independent datasets, the benchmark was engineered to preserve semantic alignment across modalities, enabling paired comparisons where differences in measured performance can be attributed primarily to response constraints rather than question intent drift. This design supports direct measurement of modality separation and exposes cases where constrained-response competence does not translate into generative adequacy under free-text requirements.

The multi-format extension contributes a reusable dataset engineering pattern for \ac{se}: conversion feasibility explicitly assessed, unsuitable items filtered to preserve construct validity, and a resulting paired dataset to supports auditable cross-modality analysis. Empirically, this instrument enabled the dissertation’s central finding that \ac{mcq} and \ac{osq} provide overlapping but non-equivalent evidence: while rank-order alignment can remain strong, absolute performance differs systematically, and the magnitude of separation varies by task category. Practically, the paired design is the enabling artifact for a task-to-modality routing map: practitioners can identify where \ac{mcq}-only evidence is likely to overstate readiness and where \ac{osq}-based qualification is required before copilots are deployed.

% The inclusion of \ac{mcq} variants supports robustness analysis by revealing sensitivity to golden answer perturbation and distractor structure to detect variability in model stability, while the addition of \acp{osq} enables evaluation of generative quality and depth. This multi-format design provides a more comprehensive evaluation instrument for different failure modes and has been explored in other domains \cite{lin_TruthfulQAMeasuringHow_2022, lin_SylinrlTruthfulQA_2026, agrawal_VQAVisualQuestion_2016, zhang_DevelopmentLargescaleMedical_2024}.

\subsection{Robustness Analysis of MCQ Evaluation Using Distractor Variation for Systems Engineering}\label{subsec:ch6_contrib2}

Building on the community, open-source \ac{mcq}-only benchmarking foundations of SysEngBench \cite{bell_IntroducingSysEngBenchNovel_2024, Bell2025SysEngBench}, this work incorporated robustness-oriented diagnostics that treat \ac{mcq} reliability as a measurable property. Specifically, the benchmark includes controlled \ac{mcq} perturbations that vary distractor structure and answer-positioning, enabling sensitivity analyses that reveal whether measured performance is stable under presentation-level changes that should be irrelevant to the underlying knowledge target.

Empirically, these diagnostics demonstrated that robustness concerns are non-trivial. A subset of evaluated models exhibited evidence patterns consistent with positional sensitivity or mixed robustness outcomes, indicating that raw \ac{mcq} accuracy can be confounded by question structural artifacts. This contribution therefore provides a practical safeguard against overconfident adoption, for systems engineering specific applications, based on brittle selection-format performance.

In deployment terms, the robustness suite turns \ac{mcq} benchmarking into a more defensible screening instrument: \ac{mcq} evidence is treated as credible only when stability under structural perturbations is demonstrated, and models exhibiting instability are either restricted to lower-risk roles or required to satisfy stronger \ac{osq}-based evidence before being used in workflows that affect configuration-managed artifacts. By examining how outcomes behave under controlled \ac{mcq} structural variation (e.g., distractor and answer-position perturbations), the benchmark functions as an instrument for detecting evaluation artifacts that could otherwise mislead practitioners into overconfident employment of language models.

\subsection{Evaluation Modality Fusion as a Basis for Model Qualification and Trustworthy Use in Systems Engineering}\label{subsec:ch6_contrib3}

This dissertation introduced evaluation modality fusion as a basis for defensible \acp{llm} qualification in systems engineering. The core conceptual contribution is the framing of \ac{mcq} and \ac{osq} as complementary sources of evidence rather than functional equivalents: because the modalities impose different response constraints, they probe different capabilities and surface different failure modes, even when questions are semantically aligned. The objective is therefore not to declare one modality ``better,'' but to interpret them jointly as qualification evidence aligned to the intended use claim and the risk posture of the target task.

Empirically, the fused perspective is justified by two observed properties of the trade space. First, \ac{osq} evaluation imposes a meaningfully higher performance bar under the same knowledge targets, producing systematic modality separation that can be statistically significant even when cross-model rank-order association remains strong. Second, the magnitude of separation is category-dependent, enabling the derivation of modality requirements that differ across task families rather than applying a single blanket rule. These findings motivate deployable decision rules. In practice, fused qualification is implemented as (i) a task-to-modality routing map that specifies the minimum defensible evaluation evidence by task class (e.g., when \ac{mcq}-only evidence is insufficient), and (ii) modality-aware evidence composition in which \ac{mcq} and \ac{osq} are combined as needed to support the intended use claim—using \ac{mcq} for scalable constrained-response measurement and robustness diagnostics, and \ac{osq} when generative adequacy and justification are central success criteria. Table~\ref{tab:task_modality_guidance} summarizes notional guidance for applying these rules by task context and oversight needs.


\subsection{Token-Efficient Qualification: Cost-Aware Response Length Trade-offs in High-Performing Models}\label{subsec:ch6_contrib4}

This dissertation evaluated response length behavior as an operational cost factor and analyzed the quality--token trade space among models evaluated under \ac{osq} with rubric-based, LLM-as-a-judge scoring. Across the evaluated models, response lengths varied substantially—from very short outputs to a design-limited 500-token cap—while more verbose responses did not consistently yield higher judged scores. This establishes token usage as a practical deployment trade space and motivates treating efficiency as a first-class selection criterion rather than an afterthought once a single ``best'' model is chosen.

The resulting contribution is a token-efficient qualification method that makes cost-aware adoption actionable: models are filtered first for acceptable \ac{osq} performance (e.g., at or above a ``Good'' or ``Pass'' threshold), and then compared on efficiency within that qualified set using the efficient frontier between judged performance and token usage (Figure~\ref{fig:pareto_quality_vs_token_usage}). Under comparable mean judged score, the most cost-effective candidates are those that achieve the required quality level using fewer median tokens per response. Practically, this yields a deployment rule aligned with performance and cost: \emph{qualify first, then select the lowest-token model among those meeting the required quality threshold}. This cost-aware framing directly supports scalable integration of copilots by minimizing avoidable compute expense and human review burden while preserving evidence-based performance requirements.


\section{Closing Thoughts}\label{sec:final_thoughts}

This dissertation supports a central conclusion: MCQ and OSQ modalities should be used together to make defensible model qualification choices for systems engineering. A statistically significant difference was observed between the two modalities, indicating that MCQ and OSQ formats are not interchangeable when accounting for a cost-efficient trade space dimension, despite strong rank-order alignment. In other words, while relative model rankings can remain consistent across formats, absolute performance differs substantially, with OSQ imposing a meaningfully higher performance bar. Treating one modality as a substitute for the other risks either overestimating competence (e.g. when MCQ performance is taken as sufficient evidence) or under-utilizing efficient screening signals (e.g. when OSQ is treated as the only credible instrument). Mixed-modality evaluation provides the stronger basis for interpreting capability and anticipating failure modes that matter for SE work products.

As SysML v2 enables \ac{mbse} models to become increasingly text-addressable and repository-managed, \ac{se} work is shifting toward an \ac{ide}-centered workflow in which architecture artifacts can be searched, diffed, tested, and refactored like other software-engineered assets. This shift preserves the foundations of systems thinking and trade-space reasoning, but elevates the importance of coding fluency, automation, and version-controlled engineering across the digital thread. The next-generation systems engineer therefore increasingly resembles a hybrid SE and software practitioner: combining architectural judgment with the tool chain competence required to scale change under configuration control.

This transition is unfolding amid accelerating operational tempos and rising system-of-systems complexity. Under these conditions, the ability to iterate quickly by exploring alternatives, refactoring architectures, and validating consequences moves from a productivity advantage to an operational necessity. Agents and copilots can help compress iteration cycles and reduce friction in routine engineering work, but only when embedded in workflows that preserve systems engineering discipline and make verification and validation integral to the loop.

Embedding AI tools such as language models into the digital thread and digital twin is therefore inherently a high-risk, high-reward proposition. When governed well, such integration can enable rapid iteration with richer context and the ability to traverse the digital thread with unprecedented efficiency.. When governed poorly, accessibility and speed without trustworthy oversight enables rapid corruption of the \ac{asot}. This dissertation was motivated by enabling the former outcome, to empower systems engineers to architect at the pace of relevancy by using benchmark-qualified copilots within workflows that emphasize governance and accountable human oversight.


\chapter*{Disclosure}\label{ch:disclosure}
\addcontentsline{toc}{chapter}{Disclosure}
% \section*{Acknowledgments}\label{sec:acknowledgments}

This work has benefited from the use of generative AI tools, including ChatGPT, Claude, and SciSpace, for brainstorming, writing assistance, and code development \cite{openai_OpenAI_2025, anthropicai_AnthropicHome_2025, scispace_SciSpace_2025}. These tools were employed to enhance conceptual clarity, improve code efficiency, and support literature synthesis. All AI-generated outputs were critically reviewed, refined, and validated to ensure accuracy and alignment with academic integrity. Their contributions were limited to supporting the research process, and final responsibility for the content, analysis, and conclusions remains with the author.
